\section*{Hoofdstuk \ref{ch:g10:lewis-and-shape} --- Lewisstructuren en de vorm van moleculen}

\begin{solution}{exo:g10:lewis-and-shape:1}
Een elektronenpaar dat door twee \omterm{def:g7:atoms-and-molecules:atom}{atomen} gedeeld wordt. Een paar
\omterm{def:g10:electron-shells:valence}{valentie-elektronen} dat bij één \omterm{def:g7:atoms-and-molecules:atom}{atoom} hoort en niet gedeeld wordt.
\end{solution}

\begin{solution}{exo:g10:lewis-and-shape:2}
H: 1 (één \omterm{def:g9:inside-the-atom:nucleus}{elektron} tekort voor zijn duet). C: 4, N: 3, O: 2 (4, 3 en 2
\omterm{def:g9:inside-the-atom:nucleus}{elektronen} tekort voor een octet).
\end{solution}

\begin{solution}{exo:g10:lewis-and-shape:3}
\chemfig{H-\charge{0=\:,90=\:,270=\:}{Cl}}: chloor heeft het gedeelde paar en drie vrije
paren, 8 \omterm{def:g9:inside-the-atom:nucleus}{elektronen}.
\end{solution}

\begin{solution}{exo:g10:lewis-and-shape:4}
$5 + 3 \times 1 = 8$ \omterm{def:g9:inside-the-atom:nucleus}{elektronen}, 4 paren (drie bindingen, één vrij paar).
\end{solution}

\begin{solution}{exo:g10:lewis-and-shape:5}
\ce{CH4}: tetraëder. \ce{NH3}: piramide. \ce{H2O}: gebogen. \ce{CO2}:
lineair.
\end{solution}

\begin{solution}{exo:g10:lewis-and-shape:6}
\chemfig{H-\charge{90=\:,270=\:}{S}-H}: net als bij water vier groepen rond zwavel,
waarvan twee vrije paren: gebogen.
\end{solution}

\begin{solution}{exo:g10:lewis-and-shape:7}
$2 \times 5 = 10$ \omterm{def:g9:inside-the-atom:nucleus}{elektronen}, 5 paren. Een enkele binding en 4 vrije paren
laten elk stikstofatoom met 6 \omterm{def:g9:inside-the-atom:nucleus}{elektronen} achter; door twee vrije paren in
gedeelde paren om te zetten, ontstaan een \omterm{def:g10:lewis-and-shape:multiple-bond}{drievoudige binding} en één vrij
paar op elk \omterm{def:g7:atoms-and-molecules:atom}{atoom}:
\chemfig{\charge{180=\:}{N}~\charge{0=\:}{N}}.
\end{solution}

\begin{solution}{exo:g10:lewis-and-shape:8}
\chemfig{H-C(-[2]H)(-[6]H)-\charge{90=\:,270=\:}{O}-H}: vier bindingen rond koolstof,
tetraëdrisch; rond zuurstof twee bindingen en twee vrije paren, gebogen.
\end{solution}

\begin{solution}{exo:g10:lewis-and-shape:9}
In \ce{CO2} heeft koolstof twee groepen (twee \omterm{def:g10:lewis-and-shape:multiple-bond}{dubbele bindingen}) en geen
vrij paar: die wijzen in tegengestelde richtingen. In \ce{H2O} heeft
zuurstof vier groepen (twee bindingen, twee vrije paren) in een
tetraëdrische schikking: de twee bindingen maken een hoek.
\end{solution}

\begin{solution}{exo:g10:lewis-and-shape:10}
\chemfig{C(-[3]H)(-[5]H)=C(-[1]H)-[7]H}: rond elk koolstofatoom drie
groepen (één \omterm{def:g10:lewis-and-shape:multiple-bond}{dubbele binding}, twee enkele bindingen): een vlakke
driehoek, hoeken van ongeveer \ang{120}.
\end{solution}

\begin{solution}{exo:g10:lewis-and-shape:11}
Koolstof in het midden, vier enkele bindingen naar chlooratomen die elk
drie vrije paren hebben: vier groepen rond koolstof, een tetraëder.
\end{solution}

\begin{solution}{exo:g10:lewis-and-shape:12}
Ethanol, \chemfig{H-C(-[2]H)(-[6]H)-C(-[2]H)(-[6]H)-\charge{90=\:,270=\:}{O}-H}, en
methoxymethaan,
\chemfig{H-C(-[2]H)(-[6]H)-\charge{90=\:,270=\:}{O}-C(-[2]H)(-[6]H)-H}.
\end{solution}

\begin{solution}{exo:g10:lewis-and-shape:13}
Een vrij paar neemt iets meer ruimte in dan een bindend paar en duwt de
bindingen dichter naar elkaar toe. Ammoniak heeft één vrij paar, water
twee: de hoek wordt kleiner van methaan naar ammoniak naar water.
\end{solution}

\begin{solution}{exo:g10:lewis-and-shape:14}
\chemfig{H-C~\charge{0=\:}{N}}: twee groepen rond koolstof, lineair.
\chemfig{H-C(-[6]H)=\charge{45=\:,315=\:}{O}}: drie groepen rond koolstof, een vlakke
driehoek.
\end{solution}

\begin{solution}{exo:g10:lewis-and-shape:15}
$3 \times 6 = 18$ \omterm{def:g9:inside-the-atom:nucleus}{elektronen}, 9 paren:
\chemfig{\charge{180=\:,270=\:}{O}=\charge{90=\:}{O}-\charge{0=\:,90=\:,270=\:}{O}}. Het centrale
zuurstofatoom heeft drie groepen (een \omterm{def:g10:lewis-and-shape:multiple-bond}{dubbele binding}, een enkele binding,
een vrij paar): het \omterm{def:g7:atoms-and-molecules:molecule}{molecuul} is gebogen.
\end{solution}

\begin{solution}{pb:g10:lewis-and-shape:1}
\textbf{1.} \chemfig{H-C(-[2]H)(-[6]H)-H}, \chemfig{H-\charge{90=\:}{N}(-[6]H)-H},
\chemfig{H-\charge{90=\:,270=\:}{O}-H}.

\textbf{2.} In elk vier groepen.

\textbf{3.} Geen bij methaan, één bij ammoniak, twee bij water.

\textbf{4.} Koolstof vormt vier bindingen; zuurstof vormt er twee (zijn
twee andere groepen zijn vrije paren, die een bouwdoos niet met
staafjes laat zien).

\textbf{5.} \chemfig{H-C(-[2]H)(-[6]H)-C(-[2]H)(-[6]H)-H}: één staafje
verbindt de twee koolstofballetjes, en er zijn zes waterstofballetjes
nodig.

\textbf{6.} Tetraëder, piramide, gebogen.

\textbf{7.} Vrije paren nemen meer ruimte in dan bindende paren en drukken
de bindingen naar elkaar toe; water, met twee vrije paren, heeft de
kleinste hoek.

\textbf{8.} De gaatjes van een koolstofballetje maken \ang{109.5}; de
echte hoek van water is \ang{104.5}: een zuurstofballetje heeft zijn eigen
twee gaatjes onder \ang{104.5} nodig.

\textbf{9.} \ang{180}: twee groepen rond koolstof wijzen in
tegengestelde richtingen, en het \omterm{def:g7:atoms-and-molecules:molecule}{molecuul} is lineair.

\textbf{10.} Twee hoekpunten die niet op dezelfde ribbe liggen, zijn
tegenoverliggende hoekpunten van één zijvlak: hun afstand is de
diagonaal van het zijvlak, $\sqrt{a^2 + a^2} = a\sqrt{2}$.

\textbf{11.} Het middelpunt ligt halverwege een lichaamsdiagonaal:
$a\sqrt{3}/2$.

\textbf{12.} H--H: $\sqrt{2} \approx \qty{1.41}{cm}$; C--H:
$\sqrt{3}/2 \approx \qty{0.87}{cm}$.

\textbf{13.} Elke waterstof zit op een hoekpunt van de kubus en de
koolstof in het middelpunt: elk hoekpunt ligt even ver van het
middelpunt.

\textbf{14.} Driehoek C\,H$_1$\,H$_2$ is gelijkbenig (C\,H$_1$ =
C\,H$_2$); de lijn van de top naar het midden van de basis staat
loodrecht op de basis.

\textbf{15.} $\sin \widehat{\text{M\,C\,H}_1} = \dfrac{\text{M\,H}_1}{\text{C\,H}_1}
= \dfrac{a\sqrt{2}/2}{a\sqrt{3}/2} = \sqrt{2/3} \approx 0.816$.

\textbf{16.} Halve hoek $\approx \ang{54.7}$; hoek H--C--H $\approx
2 \times 54.7 = \ang{109.5}$.

\textbf{17.} Het is de \ang{109.5} van de tetraëder in de tabel.

\textbf{18.} De gaatjes moeten onder \ang{109.5} ten opzichte van elkaar
geboord worden.
\end{solution}
